\section{Toroczkai's prime-gap graphs and simultaneous growth}
\label{sec:primegapgraphs}
Let $p_n$ be the $n$th prime, put $p_0=1$, and set $d_i=p_i-p_{i-1}$.
A prime-gap graph on $n\ge2$ vertices is a simple graph on $1,\ldots,n$
with $\deg(i)=d_i$. The sequence begins $1,1,2,2,4,2,4,2,4,6$.
Its degree sum is $p_n-1$. A matching consists of edges with pairwise
disjoint endpoints. Write $\nu(G)$ for its largest possible size.

\paragraph{Two named conjectures.}
Toroczkai conjectured that such a graph exists for every $n\ge2$, and
that \emph{every} such graph has a matching of size $d_{n+1}/2$.
Erd\H{o}s, Harcos, Kharel, Maga, Mezei and Toroczkai
proved both assertions under the classical Riemann hypothesis
(Theorems 1.1 and 1.6 of \cite{graphEH}).
The second implies the first. Begin with a single edge and replace
$d_{n+1}/2$ disjoint edges by paths of length two through a common new vertex.
The old degrees stay fixed and the new vertex has the required degree.
The zeta-function endpoint of the Lean repository therefore supplies the full
hypothesis of both published implications. The graph-theoretic results and
their analytic bridge are used here as theorems in print.
The RH-free bounds of Aggarwal, Frot, Gou and Wang
give existence for $n\ge\exp(\exp(30.32))$ and the matching assertion for
$n\ge\exp(\exp(34.33))$ \cite{graphAF}.

\subsection{A reserve of disjoint edges in every realization}
The same prime-gap estimate also allows many future insertions to be planned
using only the edges already present. All logarithms are natural.
Write $p=p_n$, $L=\log p$, and
\[
       N=\lceil652L^2\rceil,\qquad
       B(p)=\left\lceil\frac{p-2}{4N}\right\rceil .
\]
\paragraph{Matching and block theorem.}
Assume classical RH. Every prime-gap graph $G$ on $n\ge2$ vertices has
$\nu(G)\ge B(p)$. Whenever $p_{n+k}-p\le2B(p)$, the next $k$ vertices
can all be inserted using disjoint sets of old endpoints. No pair of new
vertices is adjacent, and every old vertex is incident to at most one
new edge. In particular, for each fixed $0<c<1/1304$ and all sufficiently
large $n$, this holds for $k=\lfloor c p/L^3\rfloor$.

\paragraph{Proof.}
For an integer $T\ge3$, let $H(T)=\sum_{i\le n,\ d_i\ge T}d_i$.
Theorem 1.8 of \cite{graphEH} bounds the sum of gaps of length at least $T$
starting at a prime in $[x,2x]$ by $163x(\log x)^2/T$ under RH,
for every real $x\ge2$. Apply this at $x=p/2^j$ for $j\ge1$ while $x\ge2$.
These intervals cover every gap start up to $p$ except possibly $1,2,3$,
whose following gaps have size at most two. Since $\sum_{j\ge1}p/2^j=p$,
\[
                  H(T)\le 163pL^2/T.                 \tag{G1}\label{eq:g1}
\]
Delete all vertices of degree at least $T$. This loses at most $H(T)$
edges, leaving maximum degree at most $T-1$. Vizing's edge-coloring
theorem partitions the remaining edges into at most $T$ matchings.
In particular
\[
 \nu(G)\ \ge\
 \left\lceil \frac{(p-1)/2-H(T)}{T}\right\rceil .    \tag{G2}\label{eq:g2}
\]
Taking $T=N$ gives $H(N)\le p/4$, hence the advertised $B(p)$.
The constant comes from the leading lower bound
$p/(2T)-163pL^2/T^2$, which is maximal at $T=652L^2$.

Partition a matching of size $(p_{n+k}-p)/2$ into pieces of respective sizes
$d_{n+1}/2,\ldots,d_{n+k}/2$. For each piece, remove its edges and join
their endpoints to the corresponding new vertex. Distinct pieces have
disjoint endpoints. This implements the claimed insertions in any order.
Finally $2B(p)\sim p/(1304L^2)$. For fixed $c<c'<1/1304$ take
$h=c'p/L^2$. Schoenfeld's RH estimate \cite[Corollary 1]{graphS}
$\pi(t)=\operatorname{li}(t)+O(\sqrt t\log t)$ gives
$\pi(p+h)-\pi(p)=(c'+o(1))p/L^3$ by subtraction at the two endpoints.
Thus the first $\lfloor c p/L^3\rfloor$ new primes lie within the allowed
distance once $p$ is sufficiently large. \hfill$\square$

\paragraph{The multiset extension question.}
Section 6 of \cite{graphDPG} asks which multisets of prescribed degrees
can be appended to a fixed initial graph by degree-preserving steps.
For a list of positive even degrees $e_1,\ldots,e_k$, the restricted
version using only original edges with pairwise disjoint endpoints has
the exact criterion $\sum_j e_j\le2\nu(G)$. Necessity follows because
all removed edges form one matching. Sufficiency is the partition
construction above. The unrestricted problem permits reusing vertices
and edges created by earlier steps. For example, starting with a single
edge one can successively add arbitrarily many vertices of degree two
by subdividing one current edge at each step. The block theorem gives
a quantitative sufficient condition for the unrestricted problem on
every prime-gap realization, together with the stronger guarantee on
which edges the extensions use.

\paragraph{Exact example.}
For $n=10{,}000$, $p=104729$ and $T=31$, the actual degree sum above
the cutoff, counting equality, is $9226$. Equation (G2) grants
$\lceil(52364-9226)/31\rceil=1392$ disjoint edges in \emph{every}
realization, without an analytic hypothesis. The next 236 primes end
at $107509$ and require $(107509-104729)/2=1390$ edges in total.
The following prime is $107563$. The accompanying Python code checks the
degree sums by sieving, constructs one realization by Havel--Hakimi,
and carries out all 236 insertions. It verifies simplicity, all $10236$
final degrees, and distinctness of the old endpoints.

\subsection{Keeping every late edge close in vertex labels}
An infinite sequence of degree-preserving steps can keep moving the same
edge endpoint indefinitely. A choice that stops moving old endpoints gives
a connected countable graph with the prescribed degrees. The following
result also controls how long each vertex has to wait for its final neighbors.

\paragraph{Stabilization theorem.}
Assume classical RH. There is a sequence of connected prime-gap graphs
$G_2,G_3,\ldots$, obtained by the insertion operation above, in which
the neighbors of every fixed vertex eventually become permanent.
Their limit is a connected, locally finite simple graph on the positive
integers with $\deg(i)=d_i$. Every limit edge $ij$ with $i<j$ satisfies
$j-i=O(j/\log j)$ as $j\to\infty$, uniformly in $i$.

\paragraph{Proof.}
Fix $A=100$ and a sufficiently large constant $P\ge e^{A}$. Use the
published all-$n$ theorem until the first stage with $p=p_n\ge P$.
Connectivity persists because each replaced edge has a new path
between its original endpoints. At that stage choose a spanning tree.
Thereafter all of its edges, and the edges of the successive trees
defined below, will be protected. Set
\[
 w(x)=\begin{cases}
 x,&0\le x<P,\\
 Ax/(A+\log(x/P)),&x\ge P,
 \end{cases}
 \qquad F(x)=x-w(x).
\]
At stage $n$ require that no edge used in the next insertion touch
any of the first $k=\pi(F(p))$ vertices. Both $w$ and $F$ are
nondecreasing, and $F$ is 1-Lipschitz. For $p\ge P$,
$w(p)\ge Ap/\log p$.

We show inductively that enough eligible edges remain. Let $u=p_k$ if
$k>0$. At stage $k$, all neighbors of the first $\pi(F(u))$ vertices
had labels at most $k$ and were permanently protected from that stage
onward by monotonicity of the cutoff. If $u<P$, this assertion is empty.
Thus at least $(v-1)/2$ edges inside the first $k$ labels still exist
at stage $n$, where $v$ is the greatest prime at most $F(u)$, using $v=1$
if there is none. This follows by summing the degrees of those earlier
vertices and allowing each internal edge to be counted twice.
The degree sum in the entire first $k$ labels is $u-1$, so the number of
edges incident to them is at most $(u-1)-(v-1)/2$.
Consequently at least $(p-2u+v)/2$ edges have both endpoints outside
the first $k$ labels.

By the RH prime-gap bound $d_{m+1}=O(\sqrt{p_m}\log p_m)$
(Theorem 2.3 of \cite{graphEH}), $v\ge F(u)-O(\sqrt p\log p)$.
As $u\le F(p)$ and $F$ is 1-Lipschitz,
$-2u+F(u)\ge-2F(p)+F(F(p))$. Hence the number of available edges
before protecting the tree or truncating high degrees is at least
\[
 \frac{p-2F(p)+F(F(p))}{2}-O(\sqrt p\log p)
 =\frac{w(p)-w(F(p))}{2}-O(\sqrt p\log p).
\]
For $k=0$ the same inequality follows directly from the total
$(p-1)/2$ edges and $0\le F(p)<2$.
For $P\le x\le p$, differentiation gives
$w'(x)=A(t-1)/t^2$ with $t=A+\log(x/P)\ge A$, hence
$w'(x)\ge w(p)/(2p)$. For $x<P$ the derivative equals one.
Integrating on $[F(p),p]$ gives
\[
                      w(p)-w(F(p))\ge \frac{w(p)^2}{2p}.
\]

The part of the protected tree induced on the $n-k$ remaining vertices
is a forest and contains at most $n-k$ edges. For all sufficiently large
$p$, uniformly in $P\le p$, the RH prime-counting estimate gives
$n-k\le3w(p)/\log p$. Indeed for $w(p)\ge p/2$ use
$\pi(p)\sim p/\log p$. Otherwise subtract the estimates for $\pi$ at
$p$ and $p-w(p)\ge p/2$ and use $w(p)\ge Ap/\log p$ to absorb the error.
Also exclude vertices of degree at least $T=\lceil(\log p)^5\rceil$.
By (G1) the resulting loss is at most $163p/(\log p)^3$ edges.
The number left is thus at least
\[
 \frac{w(p)^2}{4p}-\frac{3w(p)}{\log p}
       -O(\sqrt p\log p)-\frac{163p}{(\log p)^3}
          \gg \frac{p}{(\log p)^2}.
\]
Here $w(p)^2/(4p)\ge25w(p)/\log p$. All implied constants are
independent of $P$, so a single sufficiently large $P$ makes these
estimates valid throughout the induction. Edge coloring now gives
a matching of size $\gg p/(\log p)^7$, greater than $d_{n+1}/2$.
Make the insertion and add one edge incident to the newcomer to the
previously protected tree. The resulting trees are nested.

Since $F(p)\to\infty$, the neighbors of each fixed label eventually
become permanent. The union of the trees proves connectivity of the
limit, and stabilization preserves every finite degree. An edge $ij$,
$i<j$, is born at stage $j$, so for $p=p_{j-1}$ we have
$i>\pi(F(p))$ once $p\ge P$. With $P$ fixed,
$w(p)\sim100p/\log p$ and
$j-i\le1+\pi(p)-\pi(F(p))=O(p/(\log p)^2)=O(j/\log j)$.
\hfill$\square$

\paragraph{Which part uses RH?}
A sufficiently late starting stage for the stabilization theorem can
also be treated with older unconditional estimates. Theorems 2.2 and
2.4 of \cite{graphEH} give $d_{n+1}\le p^{5/8+o(1)}$ and
$\sum_{i\le n}d_i^2\le p^{4/3+o(1)}$. Use
$T=\lceil p^{7/20}\rceil$, so
$H(T)\le T^{-1}\sum d_i^2\le p^{59/60+o(1)}=o(p/(\log p)^2)$.
The gap estimate makes the error from rounding $F(u)$ down to a prime
$o(p/(\log p)^2)$ as well. The classical prime number theorem with
remainder $O(x/(\log x)^4)$ suffices for $n-k\le3w(p)/\log p$.
After the same exclusions, a matching has size
$\gg p^{13/20}/(\log p)^2$, exceeding $d_{n+1}/2$.
Thus any sufficiently late connected starting realization admits this
construction using classical bounds alone. The all-stage start at two
vertices comes from the published RH implication. The uniform matching
bound with $T=\lceil652L^2\rceil$ uses the RH tail estimate (G1).
